\begin{qfdefinition}[Metals markets for macro: gold and real rates, silver, platinum, copper and aluminium, with their main producers, uses and contracts as of 2013]\label{def:metals-markets}
\qffield{Statement} \emph{Gold} (XAU) \citep[chap.~10, sec.\ ``Precious Metals'', PDF pp.~319--321]{s7-gliner}: jewellery ``more than 40 percent of demand'' and investment the second driver; central banks held 31,920 tonnes in October 2013, the United States 8,134 tonnes, 73 percent of its reserves (Table~10.6); ``China has overtaken'' South Africa ``as the leading producer'', and ``China, the United States, Australia, and South Africa collectively account for more than 40 percent of global gold production'' (p.~320); India, the largest consumer, used over 500 tonnes for jewellery and 300 for investment in 2011, ``of which 50 percent is used for weddings''; gold is preferred to fiat money when real interest rates are negative (Barsky and Summers's Gibson's paradox); GLD trades at a tenth of an ounce and COMEX gold (GC) is 100 troy ounces.

\emph{Copper} (sec.\ ``Industrial Metals'', pp.~325--327): ``Dr.\ Copper'', a leading indicator; production led by Chile and Peru; China the second producer and largest importer, about 40 percent of imports; uses equipment manufacturing 52 percent, building construction 33 and infrastructure 15; LME copper in dollars a tonne, 25 tonnes a contract, and COMEX copper (HG) in cents a pound, 25,000 pounds. \emph{Aluminium} (pp.~327--328): the most abundant metal in the crust; China over 40 percent of both production and consumption and self-sufficient; energy-intensive to produce; LME aluminium 25 tonnes.

\qffield{Notation} Troy ounce 31.1 grams; LME 3-month rolling forward quoted in dollars per tonne.

\qffield{Conditions} Figures are Gliner's for 2011--2013, from the World Gold Council, GFMS, ICSG, USGS and the World Copper Bulletin as he cites them; production shares are compared with USGS in Pitfalls. \emph{Demand shares now}. Gold: the World Gold Council's full-year 2025 figures give total demand including OTC of 5,002 tonnes, investment 2,175 tonnes and central banks 863, and jewellery at a five-year low of 1,542 tonnes (India 430.5, mainland China 360.1) \citep[press release and jewellery section, 29 January 2026]{wgc-gdt-fy2025}.

Jewellery is thus about 31 percent of demand, not Gliner's ``more than 40 percent'', and investment, at about 43 percent, is now the largest use. Silver: the \emph{World Silver Survey 2026} gives 2025 total demand of 1.13 billion ounces and industrial demand of 657.4 million, about 58 percent, so most silver demand is still industrial; mine supply was 846.6 million ounces \citep[press release of 15 April 2026]{silver-institute-wss-2026}.

\emph{Platinum now}. Metals Focus's figures for the World Platinum Investment Council put 2025 total demand at 8{,}671 koz: automotive (autocatalyst) 3{,}039, jewellery 2{,}214, industrial 2{,}266 and net investment 1{,}152, and refined mine production at 5{,}557 koz, of which South Africa 3{,}957, Russia 677, Zimbabwe 515 and North America 212 \citep[Table~1, p.~5]{wpic-pq-2026q2}. Those are shares of 35.0, 25.5, 26.1 and 13.3 percent of demand and 71.2 percent of refined mine supply from South Africa; Gliner's catalytic converters at 38 percent and jewellery at 31 percent are 2013 figures on another basis. The same table forecasts a 265 koz surplus for 2026 after deficits in 2023--2025, a forecast, not data.

\emph{Copper and aluminium now}. The ICSG's \emph{World Copper Factbook 2025}, on preliminary 2024 figures, gives mine output of 23 million tonnes (Chile 24 percent), smelter output of 23.3 million (China 51 percent), refined output of 27.5 million (China 45 percent) and refined usage, by semis plants and other first users, of 27.4 million, of which China 58 percent \citep[p.~5]{icsg-factbook-2025}: Gliner's ``roughly 40 percent of imports'' for China is not comparable with a usage share.

By end use the Factbook's 2024 chart, sourced to the International Wrought Copper Council, has building construction 26 percent, infrastructure (power and telecom) 17, transport 13, industrial equipment 12, consumer and general products, cooling and electronics 23, and others 9 (p.~48); the ICA's 2023 dataset, for 2022 and the copper content of semis by end use (world total 30{,}582.6 thousand tonnes), groups its markets as Gliner does, into building construction 26.2 percent, infrastructure 17.3 and equipment manufacture 56.5 (industrial 11.5, transport 12.9 and other equipment 32.0) \citep[World sheet, rows~48--51]{ica-semis-enduse-2023}, against Gliner's 33, 15 and 52; the Factbook's 2024 chart, grouped the same way, gives 26, 17 and 57.

Gliner's year and basis are not stated, so the gap is a difference between his figures and the 2022 and 2024 data, not a measured trend. For aluminium, the IAI and Mission Possible Partnership's April 2023 transition strategy says ``Approximately 50\% of global aluminium demand today comes from the transport and construction sectors (26\% and 24\%, respectively)'', with electrical and machinery equipment 11 percent each, foil stock and packaging 8 percent each and consumer durables 6 (94 percent named); it credits the passage on p.~15 to a 2020 World Economic Forum and Accenture report and the same figures on p.~26 to the IAI Material Flow Model \citep[PDF pp.~15 and~26, endnotes 3 and~16]{iai-mpp-2023}, so they describe a model baseline of about 2020, not a dated statistical release, and give no country shares.

\emph{Contracts.} COMEX: copper (HG) 25{,}000 pounds of Grade~1 electrolytic cathode (ASTM B115), quoted in multiples of five hundredths of a cent a pound (\$12.50 a contract); silver (SI) 5{,}000 troy ounces of at least 999 fineness, \$0.005 an ounce outright (\$25) and \$0.001 for spreads; gold (GC) 100 troy ounces of at least 995 fineness, as one 100-ounce bar or three kilo bars, \$0.10 an ounce (\$10); aluminium 25 metric tons of P1020A, \$0.25 a ton (\$6.25) \citep[Rules 107102, 111101--111102, 112101--112102 and 113101--113102]{comex-metals-chapters-2026}.

LME: Copper Grade~A is refined electrolytic copper to BS EN~1978:2022 Cu-CATH-1, GB/T~467-2010 or ASTM B115-10 (2021) Grade~1, in full-plate cathodes of LME-listed brands, and High Grade Primary Aluminium is P1020A (revised January 2018) or GB/T~1196-2018 Al99.70, both in lots of 25 tonnes, 2 percent more or less \citep[Part~6, pp.~6-1 and~6-11]{lme-rulebook-2026-10}; minimum price fluctuations are ``at such levels as the Exchange may from time to time prescribe'' (Rule~11.3.7), so they are not in the rulebook.

\emph{Aluminium production by country}. USGS's \emph{Mineral Commodity Summaries} (February 2026) estimate world smelter production at 74.0 million tonnes in 2025 (72.8 in 2024), of which China 45.0 (44.0), about 61 percent in both years, India 4.2, Russia 3.9, Canada 3.3 and the United Arab Emirates 2.7; world year-end capacity 79.7 million tonnes \citep[Aluminum, p.~2]{usgs-mcs-2026-aluminum}. Gliner's ``over 40 percent'' of production is therefore well below the current share; USGS gives no consumption by country, so his consumption share is unverified. Gliner's contract sizes match these texts.

\qffield{Practice} Gold's negative relation to real rates is an empirical tendency with long breakdowns, not a fixed correlation.

\qffield{Pitfalls} Gliner writes that platinum ``is almost twice has [sic] heavy as gold'' with densities of 21 and 19 (p.~323); the ratio is about 1.11 (checks, section~14). He writes that Chile represents ``more than 40 percent of consumption'' of copper (p.~325), which cannot be right for Chile, a producer and exporter; it probably means China, which he calls the largest importer in the same paragraph, but the text does not say so; Chile at 27 percent of production and China at 43 percent of consumption were not matched to his table. ``India largest consumer (50\% used for weddings)'' compresses his statement about India's 2011 demand. ``China overtook South Africa as largest producer'' is stated by Gliner (p.~320).

\emph{Production shares against the data}. \emph{Gold}: USGS's \emph{Mineral Commodity Summaries 2013} gives 2011 and 2012 (estimated) mine output of 362 and 370 tonnes for China, 234 and 230 for the United States, 258 and 250 for Australia and 181 and 170 for South Africa, against world totals of 2,660 and 2,700 tonnes: 38.9 and 37.8 percent for the four \citep[gold, p.~2]{usgs-mcs-2013}, a little under Gliner's ``more than 40 percent'', whose source for that sentence is not stated (his Table~10.7 cites the World Gold Council, GFMS and the LBMA).

\emph{Silver}: Gliner's Table~10.9 (World Silver Survey 2013 and GFMS) is internally inconsistent. Its ounces (Mexico 162.2, China 117.0 million) sum, with Other 120.0, to its stated total of 787.0 million, on which Mexico and China are $279.2/787.0=35.5$ percent and the top five 62.6 percent; but its printed shares (27.3 and 19.7 percent, and 120.0 as 20.2 percent) imply a denominator near 594 and add to 132.5 percent; his ``47 percent'' and ``83 percent'' for the top five follow the printed shares. USGS's 2013 summary gives 2012 estimates of 4,250 tonnes for Mexico and 3,800 for China against a world total of 24,000, 33.5 percent \citep[silver, p.~2]{usgs-mcs-2013}, close to the 35.5 percent the table's own ounces imply.

Gliner's figures are quoted as he prints them, with the discrepancy noted. \emph{Current shares} (USGS, \emph{Mineral Commodity Summaries 2025}, 2024 estimates): gold, China 380 of 3,300 tonnes, and China, Russia, Australia, Canada and the United States together 41 percent; silver, Mexico 6,300 and China 3,300 of 25,000 tonnes (38 percent); copper mine output led by Chile (5,300 of 23,000 thousand tonnes) and refined output by China (12,000 of 27,000) \citep[gold, silver and copper, pp.~1--2]{usgs-mcs-2025}; these are production shares, not consumption shares.

\qffield{See also} \hyperref[def:contango-backwardation-stocks]{contango and backwardation}; \hyperref[def:tips-breakevens]{real yields}; \hyperref[def:commodity-indices-risk-return]{commodity indices}.

\qffield{Further reading} Metal inventories, exchange delivery, loco settlement and leases.

\end{qfdefinition}
